\documentclass[11pt]{article}

\usepackage[margin=1.05in]{geometry}
\usepackage{amsmath,amssymb,amsthm}
\usepackage{graphicx}
\usepackage[colorlinks=true,linkcolor=blue,citecolor=blue,urlcolor=blue]{hyperref}

\newtheorem{theorem}{Theorem}
\newtheorem{lemma}{Lemma}
\newtheorem{remark}{Remark}

\newcommand{\eps}{\varepsilon}
\newcommand{\etae}{\eta}
\newcommand{\Z}{\mathcal Z}
\newcommand{\ball}{\mathrm B}

\begin{document}

\begin{center}
{\Large Uniform almost-center domination with arbitrarily small loss}\\[0.5em]
{\normalsize Partial result packet for arXiv:1610.02557}\\[0.25em]
{\normalsize Run: \texttt{fa\_banach\_001}}
\end{center}

\section*{Source question}

The source paper is Timur Oikhberg and Pedro Tradacete,
\emph{Almost band preservers}, arXiv:1610.02557 \cite{OT}. It defines
\(\eps\)-band preserving operators and the \(\eps\)-center
\(\eps-\Z(X)\). The relevant passage is:

\begin{quote}
``In general, if \(T\in \eps-\Z(X)\), then \(T\) is \(\eps\)-BP.
We do not know whether the converse implication holds in general.''
\end{quote}

\begin{center}
\includegraphics[width=0.92\linewidth]{figures/open_problem_crop.png}
\end{center}

The source paper's Theorem 2.10 shows that, for bounded \(T\), the
\(\eps\)-BP condition implies a local domination statement: for each
\(x\in\ball(X)\) and each \(\eps'>\eps\), there is a scalar
\(\lambda_x\) such that
\[
 \big\|(|Tx|-\lambda_x |x|)_+\big\|<\eps'.
\]
The open endpoint issue is whether \(\lambda_x\) can be chosen uniformly
over \(x\in\ball(X)\) at the exact constant \(\eps\). The theorem below
obtains such a uniform choice after replacing \(\eps\) by \(\eps+\etae\).

\section*{Result}

\begin{theorem}[Uniform \(\eps+\eta\) center consequence]
\label{thm:main}
Let \(X\) be a Banach lattice, let \(T\in B(X)\), and suppose that \(T\)
is \(\eps\)-band preserving. Then for every \(\etae>0\),
\[
 T\in(\eps+\etae)-\Z(X).
\]
More quantitatively, if \(M=\|T\|>0\) and
\[
 \lambda\ge \frac{8M^2(M+1)}{\etae^2},
\]
then for every \(x\in \ball(X)\),
\[
 \big\|(|Tx|-\lambda |x|)_+\big\|\le \eps+\etae .
\]
If \(T=0\), the conclusion is trivial with \(\lambda=0\).
\end{theorem}

\section*{Proof intuition}

The proof is a quantitative version of the source paper's proof of
\cite[Theorem 2.10]{OT}. That proof shows that if a norm-one functional
\(x^*\) is disjoint from \(x\), then \(x^*\) cannot see \(Tx\) by more than
\(\eps\). We need the approximate form: if a positive norm-one functional
almost annihilates \(|x|\), say \(\langle x^*,|x|\rangle\le\alpha\), then
\(\langle x^*,|Tx|\rangle\le \eps+O(\sqrt{\alpha})\).

Now suppose \((|Tx|-\lambda |x|)_+\) has norm \(>\eps+\etae\). A positive
supporting functional for this positive part may be restricted, inside the
principal ideal generated by \(|Tx|\vee |x|\), to the region where
\(|Tx|>\lambda |x|\). On that region it almost annihilates \(|x|\), with
\(\alpha\le \|T\|/\lambda\). Choosing \(\lambda\) of order \(\etae^{-2}\)
contradicts the approximate disjoint-functional estimate.

\section*{A quantitative disjoint-functional lemma}

\begin{lemma}
\label{lem:approx-disjoint}
Let \(T\in B(X)\) be \(\eps\)-BP and put \(M=\|T\|\). Let
\(x\in\ball(X)\) and \(0\le x^*\in \ball(X^*)\). If
\(\langle x^*,|x|\rangle\le\alpha\), then for every \(s>0\),
\[
 \langle x^*,|Tx|\rangle
 \le \eps+\frac{2\alpha}{s}+M(M+1)s .
\]
\end{lemma}

\begin{proof}
Let \(u=|Tx|\), \(a=|x|\), and let \(I\) be the principal ideal generated by
\(u\vee a\). As in the proof of \cite[Theorem 2.10]{OT}, identify \(I\) with
some \(C(\Omega)\) by a lattice isomorphism \(j:C(\Omega)\to I\) satisfying
\(j1=u\vee a\). Then
\[
 \|j\|=\|u\vee a\|\le \|u\|+\|a\|\le M+1.
\]
Write \(x=j(\phi)\) and \(u=j(\phi')\). Since \(0\le u\le u\vee a\), we have
\(0\le \phi'\le1\).

Let \(\mu\) be the positive measure on \(\Omega\) induced by \(x^*\), so that
\(\int f\,d\mu=\langle x^*,j(f)\rangle\) for \(f\in C(\Omega)\). Then
\[
 \int |\phi|\,d\mu=\langle x^*,|x|\rangle\le\alpha.
\]
Choose \(h\in C(\Omega)\), \(0\le h\le1\), with \(h=1\) on
\(\{|\phi|\le s/2\}\), \(h=0\) on \(\{|\phi|\ge s\}\), and
\[
 1-h\le 2s^{-1}|\phi|.
\]
Define
\[
 \psi=(\phi_+-s1)_+-(\phi_- -s1)_+,\qquad y=j(\psi),
\]
and \(y'=j(\phi'h)\). Then \(|y|\le |x|\), hence \(\|y\|\le1\), while
\(\|x-y\|\le (M+1)s\). Also \(y\perp y'\), because \(\psi\) is supported
where \(|\phi|\ge s\) and \(h\) is supported where \(|\phi|<s\).

The lower estimate from the source proof becomes quantitative:
\[
\begin{aligned}
 \|u\wedge y'\|
 &\ge \langle x^*,y'\rangle                                      \\
 &= \int \phi' h\,d\mu                                             \\
 &= \langle x^*,u\rangle-\int \phi'(1-h)\,d\mu                     \\
 &\ge \langle x^*,u\rangle-\frac{2}{s}\int |\phi|\,d\mu             \\
 &\ge \langle x^*,u\rangle-\frac{2\alpha}{s}.
\end{aligned}
\]
Since \(T\) is \(\eps\)-BP and \(y\perp y'\), we have
\(\||Ty|\wedge y'\|\le\eps\|y\|\le\eps\). Moreover
\[
 \|\,|Tx|-|Ty|\,\|\le \|T(x-y)\|\le M(M+1)s.
\]
Thus
\[
\begin{aligned}
\eps
&\ge \||Ty|\wedge y'\|\\
&\ge \|u\wedge y'\|-\|\,|Tx|-|Ty|\,\|\\
&\ge \langle x^*,u\rangle-\frac{2\alpha}{s}-M(M+1)s.
\end{aligned}
\]
Rearranging proves the lemma.
\end{proof}

\section*{Proof of Theorem \ref{thm:main}}

Assume \(M=\|T\|>0\), fix \(\etae>0\), and choose
\[
 \lambda\ge \frac{8M^2(M+1)}{\etae^2}.
\]
Let \(x\in\ball(X)\), and set \(a=|x|\), \(u=|Tx|\), and
\[
 v=(u-\lambda a)_+.
\]
We prove \(\|v\|\le\eps+\etae\).

Suppose, toward a contradiction, that \(\|v\|>\eps+\etae\). Let \(I\) be the
principal ideal generated by \(u\vee a\), represented as \(C(\Omega)\) with
\(u\), \(a\), and \(v\) corresponding to non-negative functions
\(\phi'\), \(\phi\), and \((\phi'-\lambda\phi)_+\). Choose a positive
norm-one functional on \(X\) which sees \(v\) by more than \(\eps+\etae\),
restrict it to \(I\), and then restrict the representing measure to the open
set
\[
 A=\{\omega\in\Omega:\phi'(\omega)>\lambda\phi(\omega)\}.
\]
By the positive Hahn--Banach extension theorem for ideals, this restricted
measure extends to a positive functional \(x^*\in X^*\) with
\(\|x^*\|\le1\), and it still satisfies
\[
 \langle x^*,v\rangle>\eps+\etae.
\]
Since the measure is supported on \(A\), we have \(\phi\le \lambda^{-1}\phi'\)
on its support. Consequently
\[
 \alpha:=\langle x^*,a\rangle
 \le \lambda^{-1}\langle x^*,u\rangle
 \le \frac{M}{\lambda}.
\]
Also
\[
 \langle x^*,u\rangle\ge \langle x^*,v\rangle>\eps+\etae.
\]

Apply Lemma \ref{lem:approx-disjoint} with
\[
 s=\frac{\etae}{2M(M+1)}.
\]
The choice of \(\lambda\) gives
\[
 \frac{2\alpha}{s}\le \frac{2(M/\lambda)}{s}\le\frac{\etae}{2},
\qquad
 M(M+1)s=\frac{\etae}{2}.
\]
Therefore Lemma \ref{lem:approx-disjoint} yields
\[
 \langle x^*,u\rangle\le \eps+\etae,
\]
contradicting \(\langle x^*,u\rangle>\eps+\etae\). Hence
\[
 \big\|(|Tx|-\lambda |x|)_+\big\|\le\eps+\etae
\]
for every \(x\in\ball(X)\), as required.

\section*{Verification notes and limitations}

This is not the exact endpoint converse asked in the source paper. It proves
that the endpoint is the only possible obstruction:
\[
 \inf_{\lambda\ge0}\ \sup_{x\in\ball(X)}
 \big\|(|Tx|-\lambda |x|)_+\big\|\le \eps.
\]
What remains unknown from this argument is whether the infimum is always
attained at a finite \(\lambda\). The bound here grows like \(\etae^{-2}\),
so it gives no endpoint compactness as \(\etae\downarrow0\).

The proof was checked against the source's proof of Theorem 2.10; the only
new ingredient is the quantitative estimate in Lemma \ref{lem:approx-disjoint}.
Local run-index searches found the prior exact \(C_0(K)\)/AL-space packet for
this target but no packet for the global \(\eps+\etae\) statement.

\begin{thebibliography}{9}

\bibitem{OT}
T. Oikhberg and P. Tradacete,
\emph{Almost band preservers},
arXiv:1610.02557, 2016.
\url{https://arxiv.org/abs/1610.02557}

\end{thebibliography}

\end{document}
